\documentclass{ctexart}
\newcommand{\keywords}[1]{\par\textbf{关键词：}#1}
\begin{document}
\begin{abstract}
本文构建综合性框架，具有科学性、合理性和有效性。首先建立模型，其次进行求解，再次分析结果，最后验证模型。
\keywords{优化；模型；算法}
\end{abstract}
\section{问题重述}
首先说明问题，其次整理任务，进一步明确目标。因此可见，该问题值得研究。
\section{问题分析}
针对问题一，本文采用优化模型进行求解。首先分析数据，其次建立模型，进一步运行算法，最后得到结果。这说明该方法合理。
\section{模型假设}
首先作出合理假设，其次简化问题，进一步提高效率。这说明假设具有科学性、稳定性和普适性。
\section{符号说明}
首先定义变量，其次给出参数，最后说明单位。由此可见，符号定义较为完整。
\section{模型建立与求解}
首先构建模型，其次设置约束，再次调用算法，最后输出结果。因此，该模型能够解决问题。
\section{结果分析}
首先读取结果，其次进行比较，进一步分析原因，最后得出结论。这说明模型具有良好效果。
\section{模型检验}
首先回代结果，其次比较误差，最后验证有效性。由此说明模型有效。
\section{灵敏度与稳健性分析}
首先改变参数，其次观察结果，进一步分析变化，最后说明模型稳定。因此可见模型具有鲁棒性。
\section{模型评价与改进}
首先总结优点，其次指出不足，最后提出改进。这说明模型具有准确性、稳定性和适用性。
\begin{thebibliography}{9}\bibitem{x} 资料。\end{thebibliography}
\appendix
\section{附录}
首先给出代码，其次给出数据，最后给出结果。综上，论文完整。
\end{document}
